\section{Related Work}
\label{sec:related_work}

Vision-based traffic surveillance has evolved from isolated detection pipelines toward integrated systems that combine detection, tracking, counting, and deployment constraints. Surveys on traffic perception emphasize that real camera networks must address lighting changes, congestion, partial occlusion, mixed vehicle fleets, and large annotation requirements before they can support operational decisions reliably~\cite{Zhou2025TrafficSurveillanceSurvey,Azfar2024ComplexTrafficReview}. Early and mid-generation YOLO-based traffic systems illustrate this trajectory. Gomaa et al.~\cite{Gomaa2022} reduce counting cost in fixed-camera scenes by combining intermittent YOLOv2 detection with KLT optical flow, while Kumar et al.~\cite{Kumar2023} integrate YOLOv5 and DeepSORT for real-time traffic management. Bin Zuraimi and Kamaru Zaman~\cite{BinZuraimi2021} compare YOLOv3/v4 variants with DeepSORT and show the expected trade-off between speed and accuracy. These works support the use of single-stage detectors in traffic monitoring, but they primarily rely on horizontal boxes and do not address oriented localization or train-time class imbalance.

Deployment studies further show that accuracy alone is insufficient for municipal traffic monitoring. Balamuralidhar et al.~\cite{Balamuralidhar2021} demonstrate a real-time UAV system on embedded hardware by sharing features across detection, tracking, speed estimation, and segmentation. Alas et al.~\cite{Alas2026} extend the discussion to multi-camera urban networks, where scalable association across many CCTV streams must operate without assigning a high-end GPU to every node. These systems are important because they ground traffic vision research in hardware and cost constraints. However, their main tasks differ from the per-frame OBB detection problem studied here: they focus on tracking, counting, or re-identification, and they do not provide a segmentation-assisted augmentation strategy for rare vehicle categories in a rotated-box dataset.

Generalization across camera positions and domains remains a central limitation in traffic and vehicle detection. Koga et al.~\cite{Koga2021VehicleDomainAdaptation} study vehicle detector adaptation when source and target domains differ, using adversarial alignment to reduce prediction discrepancies. Vora et al.~\cite{Vora2023GeneralizationMultiview} show that multi-view perception systems must cope with camera placement and scene variation rather than assuming a fixed deployment geometry. These studies are not specific to Lima traffic, yet they motivate the same methodological concern: a detector trained only on observed frames can learn viewpoint-specific and background-specific regularities that do not transfer cleanly to another camera or lighting condition. For this reason, data preparation and augmentation need to increase diversity without breaking annotation geometry.

Oriented object detection addresses a different but complementary source of error. Standard horizontal boxes are convenient for generic object detection, but they enclose unnecessary background for diagonal or turning vehicles. Remote-sensing benchmarks and surveys have made this issue explicit by adopting oriented boxes for dense scenes, arbitrary object headings, and small targets~\cite{Ding2022DOTA,Wang2025OBBSurvey,Zand2022OBB}. Detector architectures such as Oriented R-CNN~\cite{Xie2021OrientedRCNN}, ReDet~\cite{Han2021ReDet}, and S2ANet~\cite{Han2022S2ANet} represent successive attempts to align proposals or features with object orientation. Loss functions based on Gaussian or Mahalanobis formulations further show that small angular errors can strongly affect high-precision rotated localization~\cite{Yang2021KLD,Wen2022MahalanobisRotatedDetection}. In traffic applications, Ahmed et al.~\cite{Ahmed2025} provide direct evidence that OBB and segmentation-assisted representations can reduce geometric measurement error compared with axis-aligned boxes.

The YOLO family is relevant because OBB accuracy must be balanced with real-time inference. YOLOv10 removes some post-processing overhead in conventional detection and strengthens the case for end-to-end single-stage models~\cite{Wang2024YOLOv10}. This paper uses the Ultralytics YOLO26s-OBB implementation for the final training family, which is consistent with the reported YOLO26 design goal of unifying real-time detection, segmentation, pose, classification, and oriented detection in one framework~\cite{Jocher2026}. Recent OBB-specific research, including task-wise sampling convolutions, also highlights the need to reduce inconsistencies between classification and localization features~\cite{Huang2025TSConv}. The remaining gap is not the absence of OBB detectors, but the lack of evidence on how these detectors behave under a highly imbalanced urban traffic taxonomy with train-only synthetic augmentation and motion-only annotation effects.

Class imbalance and small-object visibility are especially relevant for heterogeneous traffic. Espinosa et al.~\cite{Espinosa2020} review motorcycle detection in urban video and show that minority vulnerable-road-user classes require dedicated attention. Mo and Yan~\cite{Mo2020} address imbalance in oriented aerial vehicle detection using oversampling, feature amplification, and center loss, while Cheng et al.~\cite{Cheng2023SmallObjectSurvey} and Nikouei et al.~\cite{Nikouei2025SmallObjectSurvey} discuss why small objects are disproportionately lost in deep feature hierarchies. These works motivate data-level interventions, but they do not directly solve the problem of inserting additional rare vehicle instances into real traffic scenes while preserving OBB labels and the validation split.

Synthetic data augmentation offers one possible response to limited diversity, although the literature also warns that visual realism and annotation fidelity must be treated separately. Weber et al.~\cite{Weber2021ArtificialImagesVehicleDetection} show that artificial aerial vehicle images can improve detection in low-data regimes when combined with real data. Mumuni et al.~\cite{Mumuni2024} survey synthetic augmentation families and emphasize that hybrid real-plus-synthetic training can outperform either source alone when the generated samples remain useful for the downstream task. Diffusion-based surveys and controllable frameworks further argue that prompt-conditioned generation can increase semantic variability, but they also report challenges related to spatial alignment, bias, cost, and quality assessment~\cite{Alimisis2025Advances,Benkedadra2024,ZhuH2025}. These limitations matter for OBB detection because even visually plausible synthetic images can become harmful if object boundaries or rotated boxes no longer match the pixels.

Segmentation-assisted composition occupies a practical middle ground between full-image generation and standard geometric augmentation. Rahat et al.~\cite{Rahat2024} use off-the-shelf segmentation with generative background diversification to preserve foreground subjects, while Huang et al.~\cite{HuangW2025} analyze object-centric synthetic compositions and note that naive copy-paste can introduce edge artifacts and context mismatches. The implemented method is more conservative than those generative systems: it does not rely on full-frame diffusion synthesis for final training data. Instead, it uses SAM masks to extract real vehicles, inserts them into safe train-split slots, and preserves both the background outside the mask and the OBB annotation geometry. This design directly follows the literature's caution that synthetic data for detection must maintain spatial label validity.

Table~\ref{tab:related_work_summary} summarizes how the reviewed groups connect to the present study. Overall, prior work covers real-time traffic detection, edge deployment, OBB geometry, class imbalance, and synthetic data generation, but these threads are rarely combined in one reproducible pipeline for urban traffic OBB detection. The gap addressed here is therefore methodological: constructing and validating a train-only, segmentation-assisted copy-paste augmentation set for a long-tailed traffic dataset while keeping evaluation geometry and validation data unchanged.

\begin{table*}[!t]
  \caption{Synthesis of related work and remaining methodological gap.}
  \label{tab:related_work_summary}
  \centering
  \small
  \begin{tabular}{p{0.23\textwidth} p{0.35\textwidth} p{0.34\textwidth}}
    \toprule
    \textbf{Research line} & \textbf{Representative evidence} & \textbf{Gap with respect to this work} \\
    \midrule
    Traffic surveillance and tracking &
    YOLO-based counting and tracking systems demonstrate real-time feasibility in fixed and highway cameras~\cite{Gomaa2022,Kumar2023,BinZuraimi2021}. &
    Most systems use horizontal boxes and optimize counting or tracking rather than OBB label-preserving augmentation. \\
    \midrule
    Edge and multi-camera deployment &
    Embedded and multi-camera systems show that traffic perception must be computationally practical~\cite{Balamuralidhar2021,Alas2026}. &
    Deployment evidence does not resolve rare-class augmentation or moving-only annotation effects in OBB training data. \\
    \midrule
    Oriented vehicle localization &
    OBB benchmarks, detectors, and rotated losses improve localization for arbitrarily oriented objects~\cite{Ding2022DOTA,Xie2021OrientedRCNN,Han2021ReDet,Yang2021KLD,Wang2025OBBSurvey}. &
    OBB methods are usually evaluated independently of train-only synthetic composition and Lima-like traffic class imbalance. \\
    \midrule
    Synthetic and controllable augmentation &
    Artificial images, diffusion-based augmentation, and object-centric composition can increase data diversity~\cite{Weber2021ArtificialImagesVehicleDetection,Mumuni2024,Alimisis2025Advances,HuangW2025}. &
    Full-image generation and naive copy-paste may break annotation fidelity; this work keeps real masks and explicit OBB labels. \\
    \midrule
    Segmentation-assisted composition &
    SAM-based foreground extraction and object-centric pipelines motivate mask-guided placement~\cite{Rahat2024,Ahmed2025}. &
    Prior work does not provide the exact train-only SAM copy-paste protocol and validation checks used here. \\
    \bottomrule
  \end{tabular}
\end{table*}
